\section{The global frequency geometry}\label{sec:frequency}

The forcing identity determines the frequency derivative only after its
integration constant has been fixed.  The persistent circle at $h_-$ fixes
one constant; differentiable real integrals at the two infinite ends fix
the others.  We keep the positive time orientation and the lifts
\eqref{eq:frequency-lifts} throughout.

\subsection{The persistent circle and its anchor}

Use the completed-square coordinates \eqref{surf:completed-square},
and write $f_h=\Phi_h$ for that cubic. Thus
\begin{equation}\label{freq:cubic}
 Y^2=f_h(u),\qquad \omega=\frac{\dd u}{2Y}.
\end{equation}
At either critical value from \eqref{eq:critical}, direct substitution gives
\begin{equation}\label{freq:node-factorization}
 f_{h(w)}(u)=(u-a_w)^2(u-b_w),\qquad
 a_w=w^2(w-1),\quad b_w=-w^2/4.
\end{equation}
For $w=w_-$, we have $a_w<b_w<0$; hence the isolated node is separated
from the circle containing $O$ and $P$.

\begin{lemma}\label{lem:persistent}
The positive period $L$ and the one-step frequency on the persistent
circle extend real analytically across $h_-$.  The frequency restricts
to $\rho_-$ below $h_-$ and to $\rho_0$ above it.
Writing $L_0=L(h_-)$, their common endpoint values are
\begin{equation}\label{eq:center-frequency}
 L_0=\frac{2\pi}{|w_-|\sqrt{3-4w_-}},\qquad
 \rho(h_-)=\theta_T:=\frac12+
 \frac1\pi\arctan\frac1{\sqrt{3-4w_-}}\in(1/2,2/3).
\end{equation}
With a star denoting evaluation at $h_-$, the centre phase has jets
\begin{equation}\label{eq:center-phase-jets}
 \begin{aligned}
 \rho'(h_-)&=-\frac{H_*}{q_{0,*}\delta'_*L_0},
 &\beta=-2\rho'(h_-)&=\frac{2H_*}{q_{0,*}\delta'_*L_0},\\
 T=3/16:\quad \rho''(h_-)&=-\frac{3}{2q_{0,*}\delta'_*L_0},
 &\gamma=2\rho''(h_-)&=-\frac{12}{3125L_0}<0.
 \end{aligned}
\end{equation}
In particular, $\beta$ vanishes exactly at $T=3/16$.
\end{lemma}

\begin{proof}
The simple root $b_{w_-}$ extends to a real analytic root $r(h)$ near
$h_-$, with $f_h(u)=(u-r(h))g_h(u)$.  At $h_-$,
$g_h(u)=(u-a_{w_-})^2$, which is positive for $u\ge r(h)$.
This positivity persists uniformly: on a bounded interval it follows
from a positive lower bound, and at infinity from the monic quadratic
term.  A single analytic parametrization of the persistent circle is
\begin{equation}\label{freq:persistent-parametrization}
 u=r(h)+s^2,\qquad Y=s\sqrt{g_h(r(h)+s^2)},\qquad
 \omega=\frac{\dd s}{\sqrt{g_h(r(h)+s^2)}},
 \quad s\in\R\mathbf P^1.
\end{equation}
At $s=0$ the differential is regular after cancellation; in the chart
$1/s$ it is analytic and nonzero at $O$.  Thus the integrals on this
compact circle are analytic in $h$.  The marked section $P$ corresponds
to $s=p(h):=\sqrt{-r(h)}>0$.  Since increasing $s$ is positive time,
\begin{equation}\label{freq:one-step-integrals}
 L=\int_{-\infty}^{\infty}\frac{\dd s}{\sqrt{g_h(r(h)+s^2)}},\qquad
 A=\int_0^{p(h)}\frac{\dd s}{\sqrt{g_h(r(h)+s^2)}},\qquad
 \rho=\frac12+\frac{A}{L}.
\end{equation}
This fixes both the orientation and the integer lift.  At $h_-$ the
kernel is $(s^2+d_-^2)^{-1}$, where
$d_-^2=b_{w_-}-a_{w_-}=w_-^2(3-4w_-)/4$.
Consequently $L_0=\pi/d_-$ and
$A(h_-)=d_-^{-1}\arctan(\sqrt{-b_{w_-}}/d_-)$, proving
\eqref{eq:center-frequency}.  Below $h_-$, translation identifies the
two real circles, preserving $\omega$ and commuting with $+P$.
Their periods and frequencies therefore agree, so the same endpoint
jets apply to the circle that shrinks to the centre.

By \eqref{eq:wronskian}, with the single-step factor $j=1$,
\begin{equation}\label{freq:anchored-wronskian}
 J_\rho=\frac{\delta}{q_0}L^2\rho'
 =-\int_{h_-}^{h}\frac{H(s)L(s)}{q_0(s)^2}\,\dd s,
 \qquad h<h_+.
\end{equation}
Indeed, $L$ and $\rho'$ are analytic at $h_-$, $q_0(h_-)\ne0$,
and $\delta(h_-)=0$, so the integration constant is zero.
The elementary factorizations
\begin{equation}\label{freq:anchor-signs}
 H_* =w_-(2w_-+1)(4w_--3)^2,\qquad
 \delta'_* =w_-^2(4w_--3)^3<0
\end{equation}
show that $H_*$ is negative, zero, or positive according as
$T<3/16$, $T=3/16$, or $T>3/16$; equality corresponds to
$w_-=-1/2$.  Expanding \eqref{freq:anchored-wronskian} to
first order gives the first line of \eqref{eq:center-phase-jets}.
When $H_*=0$, put $t=h-h_-$.  Since $H=3t$,
$J_\rho=-3L_0t^2/(2q_{0,*}^2)+O(t^3)$, whereas
$\delta=\delta'_*t+O(t^2)$.  Comparing these analytic expansions
gives the second line; here $h_-=-2$, $q_{0,*}=-25$, and
$\delta'_*=-125/4$.
\end{proof}

The anchored identity also bounds the number of stationary energies.
Put $h_0=-32T/3$, the unique zero of $H$.  Since $q_0<0$ on
$h<h_+$, the sign of $\rho'$ is the sign of $-J_\rho$ below $h_-$
and of $J_\rho$ above it.  If $T\le3/16$, then $H<0$ below $h_-$,
so $J_\rho<0$ and $\rho_-'>0$.  If $T\ge3/16$, then $H>0$
on $(h_-,h_+)$, so $J_\rho<0$ and $\rho_0'<0$.
For $T>3/16$, we have $h_0<h_-$, $J_\rho>0$ on $[h_0,h_-)$,
and $J_\rho'>0$ on $(-\infty,h_0)$.  There is at most one zero
there, and at that zero
$\rho''=q_0J_\rho'/(\delta L^2)<0$.
For $T<3/16$, the factorization at $w_+>1$ gives $H(h_+)>0$,
so $h_-<h_0<h_+$.  Now $J_\rho>0$ on $(h_-,h_0]$ and
$J_\rho'<0$ on $(h_0,h_+)$.  Again there is at most one zero,
and it is a nondegenerate maximum.

\subsection{The middle interval and the upper return}

On $(h_-,h_+)$, the cubic $f_h$ has a single real root $r(h)<0$.
Its remaining quadratic factor is positive everywhere.  Thus
\eqref{freq:persistent-parametrization}--\eqref{freq:one-step-integrals}
hold on the whole middle interval, with $1/2<\rho_0<1$.
As $h\uparrow h_+$, write $a=a_{w_+}>0$ and $b=b_{w_+}<0$.
The simple root tends to $b$ and $g_h(u)$ tends to $(u-a)^2$.
The short arc $0\le s\le\sqrt{-r(h)}$ stays a fixed distance from
$u=a$; hence its integral has the finite limit
\[
 A(h)\longrightarrow\int_0^{\sqrt{-b}}\frac{\dd s}{a-b-s^2}.
\]
The full-period kernel tends to $|s^2-(a-b)|^{-1}$.
On a fixed interval around $\sqrt{a-b}$, Fatou's lemma forces
$L\to\infty$, since the limiting nonnegative kernel is not integrable.
This holds along every approaching sequence.  Therefore
$\rho_0\to1/2$.  For $T<3/16$, \eqref{eq:center-phase-jets} gives
$\rho_0'(h_-)>0$, while the upper endpoint $1/2$ lies below
$\theta_T$.  The mean value theorem supplies a point with negative
derivative.  The preceding one-zero bound now proves that the middle
maximum exists and is unique.

On the upper interval the ordered roots satisfy $r_1<0<r_2<r_3$.
Parametrize its identity circle by
\[
 u=r_3+s^2,\qquad
 Y=s\sqrt{(s^2+r_3-r_1)(s^2+r_3-r_2)}.
\]
Let $\psi_h(s)$ denote the positive reciprocal square root in this
formula.  Since $2P$ belongs to this component, $T\ge r_3>r_2$.
The point $2P=(T,T(h-1)/2)$ has the analytic signed parameter
\begin{equation}\label{freq:upper-real-arc}
 s_2(h)=\frac{T(h-1)}{2\sqrt{(T-r_1)(T-r_2)}},\qquad
 s_2(h)^2=T-r_3,\qquad
 \rho_+(h)=\frac{\displaystyle\int_{-\infty}^{s_2(h)}\psi_h(s)\,\dd s}
 {\displaystyle\int_{-\infty}^{\infty}\psi_h(s)\,\dd s}.
\end{equation}
The denominator defining $s_2$ is strictly positive.  Thus the formula
has no spurious singularity at $h=1$, where evenness gives
$\rho_+(1)=1/2$.  Compact-circle regularity, including the $1/s$ chart,
proves that $L$ and $\rho_+$ are analytic for every $h>h_+$, and
$0<\rho_+<1$.  This is the real $2P$ arc, not an arc from $O$ to the
other component.

At $h\downarrow h_+$, the roots tend to $b,a,a$ and
$s_2\to-\sqrt{T-a}<0$, since $T=w_+a>a$.
The full-period kernel tends to
$1/(|s|\sqrt{s^2+a-b})$, so Fatou's lemma again gives $L\to\infty$.
The arc in the numerator ends a fixed negative distance from zero.
On bounded subarcs its kernel is uniformly bounded, and at negative
infinity it is bounded by $C/s^2$.  Dominated convergence gives a
finite positive numerator limit.  Hence $\rho_+\to0$.

The double-step identity \eqref{eq:wronskian} gives
$J_{\rho_+}'=-2HL/q_0^2<0$ wherever $q_0\ne0$; here $H>0$ on
the entire upper interval.  The apparent singularity $h_q=9/8$
does not belong to the discriminant.  Indeed,
$K_{\mathrm{tw}}:=\delta L^2\rho_+'$ is analytic there, and
\begin{equation}\label{freq:apparent-point}
 q_0K_{\mathrm{tw}}'-8K_{\mathrm{tw}}=-2HL,
 \qquad
 \rho_+'(h_q)=\frac{H(h_q)}{4\delta(h_q)L(h_q)}>0.
\end{equation}
The second equality follows by analytic continuation and evaluation
at $h_q$.  Thus $J_{\rho_+}=K_{\mathrm{tw}}/q_0$ tends to
$-\infty$ from the left and $+\infty$ from the right.
On $(h_+,h_q)$, a point with $J_{\rho_+}\ge0$ would, by strict
decrease, give $J_{\rho_+}>0$ at all earlier points and hence
$\rho_+'<0$ there.  This contradicts $\rho_+>0$ and its lower
endpoint zero.  Therefore $\rho_+'>0$ up to and including $h_q$.
Above $h_q$, $J_{\rho_+}$ strictly decreases from $+\infty$ and has
at most one zero.  At such a zero,
$\rho_+''=q_0J_{\rho_+}'/(\delta L^2)<0$.

\subsection{Differentiable infinity constants}

The remaining existence decisions require constants, not just the
sign of the forcing.  Although the observables have compact support,
these limits are needed to classify all finite stationary circles.

\begin{lemma}\label{lem:infinity}
For each fixed $T>0$, the two weighted derivatives satisfy
\begin{equation}\label{eq:infinity-constants}
 \begin{aligned}
 J_{\rho_-}(h)&=\frac{\log T}{8}
       +O_T\!\left(\frac{\log|h|}{|h|}\right)&& (h\to-\infty),\\
 J_{\rho_+}(h)&=-\frac{\log T}{4}
       +O_T\!\left(\frac{\log h}{h}\right)&& (h\to+\infty).
 \end{aligned}
\end{equation}
The limits hold also at $T=1$.
\end{lemma}

\begin{proof}
Write $h=\nu/\xi$, where $\nu\in\{-1,1\}$ and $\xi>0$ tends to
zero, and put $\ell=\log(1/\xi)$, $c=\log T$.
The notation $R=O_{C^1,T}(\xi^a\ell^b)$ means
$|R|+|\xi\partial_\xi R|\le C_T\xi^a\ell^b$.
Substitution of $u=\nu T\xi+\xi^2z$ into $\xi^{-2}f_{\nu/\xi}(u)$
gives the exact polynomial
\[
 \Phi_\nu(\xi,z)=\frac{z^2}{4}-T^3
 +\nu\xi(T^3-2T^2z)+\xi^2(3T^2z-Tz^2)
 +3\nu T\xi^3z^2+\xi^4z^3.
\]
Its roots $z=\pm2T^{3/2}$ at $\xi=0$ are simple.  Their analytic
continuations $z_\pm(\xi)$ therefore give
\[
 r_2=\nu T\xi+\xi^2z_-(\xi),\quad
 r_3=\nu T\xi+\xi^2z_+(\xi),\quad
 r_1=T-\frac1{4\xi^2}-r_2-r_3.
\]
These are all three ordered real roots for sufficiently small $\xi$.
Analytic Taylor remainders yield
\begin{equation}\label{freq:root-scales}
 \begin{aligned}
 D_r:=r_3-r_1&=\frac{1+O_{C^1,T}(\xi^2)}{4\xi^2},&
 d_r:=r_3-r_2&=4T^{3/2}\xi^2(1+O_{C^1,T}(\xi)),\\
 \kappa_{\mathrm{mod}}:=\sqrt{d_r/D_r}
 &=4T^{3/4}\xi^2(1+O_{C^1,T}(\xi)).
 \end{aligned}
\end{equation}
In particular $\xi\kappa_{\mathrm{mod}}'/\kappa_{\mathrm{mod}}
=2+O_T(\xi)$.  The derivative information comes from the root
equation, not from differentiating an uncontrolled asymptotic.

For $\lambda>0$ define
\[
 K_c(\lambda)=\int_0^\infty
 \frac{\dd t}{\sqrt{(1+t^2)(\lambda^2+t^2)}}.
\]
The substitution $u=r_3+D_rt^2$ gives
$L=2K_c(\kappa_{\mathrm{mod}})/\sqrt{D_r}$.
We need the constant and a derivative estimate:
\begin{equation}\label{freq:complete-integral}
 K_c(\lambda)=\log(4/\lambda)+E(\lambda),\qquad
 |E|+|\lambda E'|\le C\lambda^2\log(1/\lambda).
\end{equation}
To prove it, subtract $\operatorname{arsinh}(1/\lambda)$ and split
at $t=1$.  The remainder is
\[
 R(\lambda)=\int_0^1\frac{a(t)\,\dd t}{\sqrt{t^2+\lambda^2}}
 +\int_1^\infty\frac{\dd t}{\sqrt{1+t^2}\sqrt{t^2+\lambda^2}},
 \qquad a(t)=(1+t^2)^{-1/2}-1.
\]
Since $|a(t)|\le Ct^2$, setting $\lambda=0$ is legitimate.
The primitive
$\log(t/(1+\sqrt{1+t^2}))$ of $1/(t\sqrt{1+t^2})$
shows that $R(0)=\log2$.
For $t<\lambda$, bound the difference of the reciprocal square roots
by $1/t$; for $\lambda<t<1$, bound it by $\lambda^2/(2t^3)$.
After multiplication by $|a(t)|$, these give
$|R(\lambda)-R(0)|\le C\lambda^2\log(1/\lambda)$; the tail is
$O(\lambda^2)$.  Differentiation is legitimate for each positive
$\lambda$.  The first derivative kernel, after multiplication by
$\lambda$, is bounded by
$\lambda^2|a(t)|/(t^2+\lambda^2)^{3/2}$.
The same split gives respectively $O(\lambda^2)$ and
$O(\lambda^2\log(1/\lambda))$, with an $O(\lambda^2)$ tail.
Finally $\operatorname{arsinh}(1/\lambda)-\log(2/\lambda)$ and its
$\lambda\partial_\lambda$ derivative are $O(\lambda^2)$.
This proves \eqref{freq:complete-integral}, and hence
\begin{equation}\label{freq:period-C1}
 L=\xi\bigl(8\ell-3c+R_L(\xi)\bigr),\qquad
 R_L=O_{C^1,T}(\xi).
\end{equation}

The two incomplete arcs have different moving scales.  Set $U_-=0$
and $U_+=T$, so the relevant marked point is $P$ at the negative
end and $2P$ at the positive end.  There its $Y$ coordinate is
positive, and $B_\nu:=U_\nu-r_3>0$.  The exact positive integrals are
\begin{equation}\label{freq:infinity-arcs}
 \begin{aligned}
 A_\nu&=\frac12\int_{r_3}^{U_\nu}\frac{\dd u}{\sqrt{f_h(u)}}
 =\frac1{2\sqrt{D_r}}\int_0^{B_\nu}
 \frac{(1+t/D_r)^{-1/2}}{\sqrt{t(t+d_r)}}\,\dd t,\\
 \rho_\nu&=\frac12+\frac{A_\nu}{L}.
 \end{aligned}
\end{equation}
Put $\eta=d_r/B_\nu$ and $\zeta=B_\nu/D_r$.
The integral without the factor $(1+t/D_r)^{-1/2}$ is
$2\operatorname{arsinh}(\eta^{-1/2})$.  After $t=B_\nu x$, its
error is
\[
 \mathcal E(\eta,\zeta)=\int_0^1
 \frac{(1+\zeta x)^{-1/2}-1}{\sqrt{x(x+\eta)}}\,\dd x,
 \qquad
 |\mathcal E|+|\eta\partial_\eta\mathcal E|
       +|\zeta\partial_\zeta\mathcal E|\le C\zeta.
\]
Indeed, the numerator and its $\zeta\partial_\zeta$ derivative
are bounded by $C\zeta x$; the $\eta$ derivative adds a factor
bounded by $\eta/(2(x+\eta))\le1/2$.
The resulting integrals are dominated by $C\zeta$.
The analytic roots give
\[
 \begin{array}{ll}
 B_-=T\xi(1+O_{C^1,T}(\xi)),&
 \eta_-=4\sqrt T\,\xi(1+O_{C^1,T}(\xi)),\quad
 \zeta_-=O_{C^1,T}(\xi^3),\\
 B_+=T(1+O_{C^1,T}(\xi)),&
 \eta_+=4\sqrt T\,\xi^2(1+O_{C^1,T}(\xi)),\quad
 \zeta_+=O_{C^1,T}(\xi^2).
 \end{array}
\]
All these positive parameters have bounded $\xi$ logarithmic
derivatives.  Thus the displayed two-parameter bound remains valid
with first-derivative control along both paths.
Using
\[
 \operatorname{arsinh}(\eta^{-1/2})=
 \log2-\tfrac12\log\eta+
 \log\frac{1+\sqrt{1+\eta}}2,
\]
whose last term and $\eta\partial_\eta$ derivative are $O(\eta)$,
we obtain
\begin{equation}\label{freq:arcs-C1}
 A_-=\xi(\ell-c/2+R_-),\qquad
 A_+=\xi(2\ell-c/2+R_+),\qquad
 R_\pm=O_{C^1,T}(\xi).
\end{equation}

We can now differentiate before taking a limit.  Let $a_-=1$,
$a_+=2$, $B=8\ell-3c$ and $C_\nu=a_\nu\ell-c/2$.
Equations \eqref{freq:period-C1} and \eqref{freq:arcs-C1} imply
\[
 L\partial_\xi A_\nu-A_\nu\partial_\xi L
 =\xi^2(B\partial_\xi C_\nu-C_\nu\partial_\xi B)
       +O_T(\xi^2\ell)
 =\xi(3a_\nu-4)c+O_T(\xi^2\ell).
\]
Every remainder term here uses the proved bounds
$|R_L|+|R_\nu|=O_T(\xi)$ and
$|R_L'|+|R_\nu'|=O_T(1)$.
Since $\partial_h=-\nu\xi^2\partial_\xi$ and
$\delta/q_0=\nu(1+O_T(\xi))/(8\xi^3)$,
\[
 L^2\rho_\nu'=-\nu(3a_\nu-4)c\xi^3+O_T(\xi^4\ell),
 \qquad
 J_{\rho_\nu}=-\frac{(3a_\nu-4)c}{8}+O_T(\xi\ell).
\]
This proves \eqref{eq:infinity-constants} without division by $c$,
so it includes $T=1$.  The same integral ratios give
$\rho_-\to5/8$ and $\rho_+\to3/4$; no derivative is taken of these
last ratio limits.
\end{proof}

\subsection{Exactly the stationary circles}

\begin{proposition}\label{prop:frequency}
The three return-frequency lifts have exactly the derivative signs
in Table~\ref{freq:classification-table}.  Every stationary point
listed there is a nondegenerate maximum, and there are no other
stationary points.  At $T=3/16$, the persistent circle at $h_-=-2$
is an additional nondegenerate stationary circle, although the whole
fibre at that energy is singular.
\end{proposition}

\begin{proof}
Only the two outer existence decisions remain.  If $T>3/16$, the
anchored identity gives $J_{\rho_-}>0$ on $[h_0,h_-)$ and strict
increase on $(-\infty,h_0)$.  For $3/16<T<1$ its limit
$(\log T)/8$ is negative, so it crosses zero exactly once.
For $T\ge1$ the limit is nonnegative, and strict increase makes
$J_{\rho_-}>0$ throughout the remaining interval.  Since
$q_0/(\delta L^2)<0$, these are respectively the lower maximum
and strict decrease.  In particular, the zero limit at $T=1$
does not produce a finite zero.

Above $h_q$, $J_{\rho_+}$ strictly decreases from $+\infty$ to
$-(\log T)/4$.  For $T\le1$ it is strictly positive at every finite
point, while for $T>1$ it crosses zero exactly once.  The signs
below $h_q$ and the nondegeneracy at any zero were proved above.
The middle existence argument and the anchored lower signs give
the other entries.  All inequalities are strict away from the
specified zeros, excluding horizontal inflection points.
Finally, \eqref{eq:center-phase-jets} proves the assertion at
$T=3/16$ on the analytic persistent-circle family.
\end{proof}

\begin{table}[htbp]
\centering
\begin{tabular}{c c c c}
\toprule
$T$ & $h<h_-$: $\rho_-$ & $h_-<h<h_+$: $\rho_0$
    & $h>h_+$: $\rho_+$\\
\midrule
$0<T<3/16$ & $\uparrow$ & $\uparrow\downarrow$ & $\uparrow$\\
$T=3/16$ & $\uparrow$ & $\downarrow$ & $\uparrow$\\
$3/16<T<1$ & $\uparrow\downarrow$ & $\downarrow$ & $\uparrow$\\
$T=1$ & $\downarrow$ & $\downarrow$ & $\uparrow$\\
$T>1$ & $\downarrow$ & $\downarrow$ & $\uparrow\downarrow$\\
\bottomrule
\end{tabular}
\caption{Strict frequency classification on the three regular energy
intervals.  The lower and middle lifts describe $F_T$; the upper lift
describes its first circle return $F_T^2$.  An arrow denotes an
everywhere strict derivative sign, and $\uparrow\downarrow$ denotes
exactly one nondegenerate maximum.  The persistent stationary circle
at $T=3/16$, $h_-=-2$, lies outside these open intervals.}
\label{freq:classification-table}
\end{table}

For the phase $\sigma$ of $F_T^2$, the lower and middle derivatives
are multiplied by two, while the upper phase is already $\rho_+$.
Consequently $\sigma''<0$ at every stationary circle.  In the five
parameter rows the stationary-circle set $\mathcal S_T$ consists,
respectively, of one middle circle, the persistent circle at $h_-$,
two lower circles at the same energy, no circle, and two upper
circles at the same energy.  These are exactly the circles in the
stationary-phase sum of Theorem~\ref{thm:main}.  The shrinking circle
at $T=3/16$ instead supplies the endpoint jet in
\eqref{eq:center-phase-jets}; it must not be identified with the
persistent circle on that energy level.
